% SHARD-ID: RC-03H-CCM-TENSOR-KERNEL
% SHARD-TITLE: CCM in Tensor-Network Terms II: Kernels and the Reconstructed Bond
% SHARD-SUMMARY: Identifies the exact finite Weil kernel with polynomial relations on distinct spectral support and gives the under-resolved, critical, and over-resolved regimes.
% SHARD-SUMMARY: Separates moment multiplicities, Jordan data, physical letters, and graded net weights, retaining explicit counterexamples to stronger reconstruction claims.
% SHARD-KEYWORDS: CCM, kernel, minimal polynomial, GNS, quotient, moment problem, grading, Jordan blocks
% SHARD-NOTES: none
% SHARD-SCRIPTS: scripts/ccm_tensor_network.py

\section{CCM in Tensor-Network Terms II: Kernels and Reconstruction}
\label{sec:ccm-tn-kernel}

The finite kernel is an annihilating filter on ring lengths. Its
meaning follows from evaluation on the retained support, before any
claim that it recovers the physical tensor. Use
\cref{def:ccm-tn-moments} with $r$ distinct circle points and positive
weights $m_a$.

\begin{theorem}[The three finite window regimes]
\label{thm:ccm-tn-atomic-reconstruction}
\claimstatus{thm:ccm-tn-atomic-reconstruction}{proved}
For $K$ lengths $0,\ldots,K-1$,
\[
 c^*\TNGram c=\sum_{a=1}^r m_a|p_c(z_a)|^2,
 \quad \operatorname{rank}\TNGram=\min(K,r),
 \quad \ker\TNGram=\{c:p_c(z_a)=0\ \forall a\}.
\]
At $K=r+1$, the radical is the line of coefficients of $p_*$.
For $K\leq r$ it is zero. For $K>r+1$ it consists of the multiples
$p_*s$ with $\deg s\leq K-1-r$; an arbitrary nullvector may have
additional arbitrary roots. The moments and distinct nodes uniquely
determine the weights through a Vandermonde system.
\end{theorem}
\begin{proof}
\begin{enumerate}
\item[$\langle1\rangle$1.] Substitute
$t_{k-j}=\sum_a m_a\bar z_a^{\,j}z_a^k$ and expand the quadratic
form. The evaluation matrix $A_{aj}=z_a^j$ gives
$W^{(K)}=A^*\operatorname{diag}(m_a)A$.
\item[$\langle1\rangle$2.] Distinct nodes give Vandermonde rank
$\min(K,r)$. Positive weights give exactly the evaluation kernel.
\item[$\langle1\rangle$3.] Polynomial division shows that all
vanishing polynomials are multiples of $p_*$. The degree bound yields
the three cases. The first $r$ equations
$t_k=\sum_a m_a z_a^k$, $k=0,\ldots,r-1$, have invertible
Vandermonde coefficient matrix, proving uniqueness of the weights.
\end{enumerate}
\end{proof}

Below the critical window a unit least eigenvector minimises
$\sum_a m_a|p_c(z_a)|^2$, a residual on the unknown spectral support.
Subtracting the least eigenvalue creates a different positive Toeplitz
matrix. Its radical belongs to that fitted moment problem. When its
radical is one-dimensional, the usual circle-root conclusion applies;
neither exact recovery of the original support nor simplicity follows
from the subtraction alone.

\begin{proposition}[The moment shift and the information it retains]
\label{prop:ccm-tn-trace-semisimplification}
\claimstatus{prop:ccm-tn-trace-semisimplification}{proved}
Multiplication by $z$ on $\TNmoment$ is unitary, has cyclic vector
$[1]$, and has each distinct $z_a$ once. Its cyclic matrix elements
recover $t_k$, while its unweighted trace generally does not.
For a finite retained operator with this spectrum,
$p_*(E)=0$ iff $E$ is semisimple. Recovered positive integer weights
determine the semisimple similarity class
$\bigoplus_a z_a I_{m_a}$, not Jordan block sizes or native eigenvectors.
\end{proposition}
\begin{proof}
\begin{enumerate}
\item[$\langle1\rangle$1.] Evaluation identifies the quotient with
$\mathbb C^r$ in the weighted product of
\cref{def:ccm-tn-moments}. Multiplication is diagonal by $z_a$,
so it is unitary. Polynomials interpolate arbitrary values; $[1]$ is
cyclic and $\langle[1],z^k[1]\rangle=\sum_a m_a z_a^k$.
\item[$\langle1\rangle$2.] A Jordan block at $z_a$ contributes its
size times $z_a^k$ to every power trace. Its largest size supplies
the exponent at $z_a$ in the minimal polynomial. The square-free
$p_*$ annihilates $E$ exactly when all these sizes are one.
\item[$\langle1\rangle$3.] Step 3 of
\cref{thm:ccm-tn-atomic-reconstruction} recovers the weights. Together
with the nodes they specify the displayed semisimple similarity class.
\end{enumerate}
\end{proof}

For $E=\left(\begin{smallmatrix}1&1\\0&1\end{smallmatrix}\right)$,
all trace moments are $2$. The critical kernel is $z-1$, but
$E-I\ne0$; the actual minimal polynomial is $(z-1)^2$.
This is the existing Jordan obstruction of
\cref{prop:weil-blind-jordan}, now expressed as a failed operator
annihilation. With semisimplicity the map $[p]\mapsto p(E)$ does
identify $\TNmoment$ with the cyclic operator algebra generated by
$E$. It identifies that algebra with the whole retained vector space
only in the simple-spectrum cyclic case.

\subsection{Graded networks and the sign of the form}

\begin{proposition}[Positivity tests the net graded weights]
\label{prop:ccm-tn-graded-net-criterion}
\claimstatus{prop:ccm-tn-graded-net-criterion}{proved}
For circle-supported even and odd sectors with distinct union support
$\{z_a\}_{a=1}^r$, put $d_a=m_0(z_a)-m_1(z_a)$.
Their supertrace form is
\[
 c^*W_{\mathrm{str}}^{(K)}c=\sum_a d_a|p_c(z_a)|^2.
\]
It is positive for all windows iff every $d_a\geq0$; already $K\geq r$
detects this condition. Sector unitarity and nonnegative individual
ring counts do not imply that positivity. The net moments do not
determine the two sector multiplicities separately.
\end{proposition}
\begin{proof}
\begin{enumerate}
\item[$\langle1\rangle$1.] Subtract the sector evaluation formulas.
If all $d_a\geq0$, every quadratic form is positive. At $K\geq r$,
Lagrange interpolation isolates each support point, so positivity
requires each $d_a\geq0$.
\item[$\langle1\rangle$2.] A physical qubit with parity
$P=\operatorname{diag}(1,-1)$ and the single even letter $A=P$
has doubled transfer $\operatorname{Ad}(P)$. Its even block is $I_2$
and its odd block is $-I_2$. The parity-closed rings are
$2-2(-1)^n\geq0$, yet
$W_{\mathrm{str}}^{(2)}=\left(\begin{smallmatrix}0&4\\4&0\end{smallmatrix}\right)$
has eigenvalues $4,-4$.
\item[$\langle1\rangle$3.] The sector multiplicities $(2,1)$ and
$(1,0)$ at the same point both give net weight $1$.
\end{enumerate}
\end{proof}

For curve counts, the known even term is moved across to give an
ordinary odd-sector trace with positive weights. That rearrangement
is substantive data. For a mixed divisor it cannot be replaced by
the assertion ``the odd block is unitary''. In the Riemann setting
the geometric pole and archimedean terms have not thereby become
specified physical bond sectors.

\begin{observation}[The physical tensor is not identified by ring norms]
\label{obs:ccm-tn-no-letter-reconstruction}
\claimstatus{obs:ccm-tn-no-letter-reconstruction}{proved}
Even full exact ring data do not determine labelled physical MPS
letters. On a one-dimensional auxiliary space, the two-letter tensors
$(1,0)$ and $(1/\sqrt2,1/\sqrt2)$ both have transfer $E=1$ and all
ring norms $1$, but give different product-state amplitudes in the
fixed physical basis. Thus reconstructing a weighted cyclic support
model is not recovery of those physical letters.
\end{observation}
\begin{proof}
Both transfers are the scalar sum of squared absolute values, namely
$1$. The states are respectively $|0\rangle^{\otimes n}$ and
$((|0\rangle+|1\rangle)/\sqrt2)^{\otimes n}$.
\end{proof}

There is also a useful positive-model caution: an interior mode
$|z|<1$ has a Poisson measure on the unit circle, whose multiplication
operator is unitary and reproduces the same moments. This is a
representation of the moments, not proof that the original mode was
on the circle. It illustrates why the duality input and the chosen
space cannot be suppressed in the tensor-network dictionary.

\subsection{Rejected stronger formulations retained in the ledger}

\begin{observation}[Rejected: an unqualified operator annihilator]
\label{obs:ccm-tn-unqualified-kernel}
\claimstatus{obs:ccm-tn-unqualified-kernel}{refuted}
Rejected assertion: the critical trace-moment null polynomial always
annihilates the original retained operator. The Jordan example above
has null polynomial $z-1$ but $E-I\ne0$. The surviving statement is
\cref{prop:ccm-tn-trace-semisimplification}, with semisimplicity.
\end{observation}

\begin{observation}[Rejected: sector unitarity implies graded positivity]
\label{obs:ccm-tn-positive-supertrace}
\claimstatus{obs:ccm-tn-positive-supertrace}{refuted}
Rejected assertion: unitary even and odd transfer blocks imply a
positive supertrace Weil form. The physical letter $A=P$ above has
unitary sector blocks and the indefinite two-length matrix
$\left(\begin{smallmatrix}0&4\\4&0\end{smallmatrix}\right)$.
The replacement is \cref{prop:ccm-tn-graded-net-criterion}.
\end{observation}
